\documentstyle[%


righttag%


,11pt%


,amssymb%


,russian%               remove in Eng.


,izvru%                 Set izven% in Eng.


]{amsart}





%





\newcommand{\sgn}{\operatorname{sgn}}


\newcommand{\mes}{\operatorname{mes}}


\newcommand{\im}{\operatorname{Im}}


\newcommand{\re}{\operatorname{Re}}


\newcommand{\tts}[1]{}





%\hoffset=-15mm%                  For Eng.


\setcounter{page}{49}


\god{2003}


\nomer{3~(490)} %                        Remove in Eng.


%\nomer{Vol.\,47, No.\,3}%               Set in Eng.


%\str{\pageref{begin}--\pageref{end}}%  Set in Eng.


\udk{517.95}





\author{\it К.Б.\,Сабитов, Н.Г.\,Шмел\"ева}


% NB:   ^^ remove in Eng.


\title{Краевые задачи для уравнения 


Лаврентьева--Бицадзе\\ с комплексным параметром}


\thanks{Работа выполнена при финансовой поддержке Российского фонда


фундаментальных исследований, грант \No\,02-01-97901.}





\date{}


%\pagestyle{myheadings}%         Set in Eng.


\pagestyle{plain} %              Remove in Eng.


\begin{document}


%\thispagestyle{plain}%          Set in Eng.


\maketitle


\label{begin}








В ([1], c.\,69) для области $G\subset{\mathrm R}^2$, звездной относительно 


начала координат, получена формула


\begin{equation}


u(x,y)=u_0(x,y)-\int_0^1 u_0(xt,yt)\frac {\partial }


{\partial t}J_0[\sqrt {\lambda (x^2+y^2)(1-t)}]dt,


\end{equation}


связывающая все регулярные дважды (непрерывно дифференцируемые) в $G$


решения уравнения Гельмгольца


\begin{equation}


u_{xx}+u_{yy}+\lambda u=0,


\end{equation}


где $\lambda$ --- числовой параметр с гармоническими в $G$ функциями


$u_0(x,y)$.





В [2] показано, что имеет место аналог формулы (1) для


телеграфного уравнения


\begin{equation}


u_{xx}-u_{yy}+\lambda u=0


\end{equation}


и для уравнения Лаврентьева--Бицадзе с числовым параметром


\begin{equation}


Lu\equiv u_{xx}+\sgn y\cdot u_{yy}+\lambda u=0.


\end{equation}


Там же получена теорема единственности 


решения задачи Трикоми для уравнения (4) при 


$\lambda <0$ и указана идея свед\'ения этой задачи 


к задаче Трикоми для уравнения Лаврентьева--Бицадзе.





В данной работе проверена справедливость интегрального представления решений


уравнения (4) с комплексным параметром


$\lambda $ и доказана его обратимость. 


Для комплексных $\lambda$ получена теорема единственности


решения задачи Трикоми и 


доказана теорема существования решения задачи Трикоми


для уравнения (4) при более слабых ограничениях на граничные данные.


Также указаны приложения интегрального представления 


решений уравнения (4) при


решении задачи Франкля и обобщенной задачи Трикоми для этого уравнения.


\medskip





{\bf 1.} Рассмотрим уравнение (4) в области $D$,


ограниченной ляпуновской кривой $\Gamma$, лежащей в полуплоскости $y>0$ с


концами в точках $A{=}(0,0)$ и $B{=}(1,0)$ и характеристиками $AC(x+y=0)$


и $CB(x-y=1)$ уравнения (4) при $y<0$.


Пусть область $D_+ = D\cap \{y>0\}$, $D_- =D\cap \{y<0\}$ и $D_+$


является звездной относительно начала координат.





\begin{lem}


Если функция $u_0(x,y)$ является регулярным


в $D_-$ решением уравнения струны


\begin{equation}


u_{0xx}-u_{0yy}=0,


\end{equation}


то функция вида


\begin{equation}


u(x,y)=u_0(x,y)-\int_0^1 u_0(xt,yt)\frac {\partial }


{\partial t}J_0[\sqrt {\lambda (x^2-y^2)(1-t)}]dt


\end{equation}


является регулярным в $D_-$ решением телеграфного уравнения $(3)$.


\end{lem}





{\bf Доказательство.} B целях удобства проверки формулы (6) положим


$u_0(x,y)=\varphi (x+y)+\psi (x-y)$, где


$\varphi ,\psi \in C^2$, и перейдем в характеристические


координаты $\xi =x+y$ и $\eta =x-y$. Тогда получим


\begin{multline*}


u(x,y)=u\left (\frac {\xi +\eta }{2},\frac {\xi -\eta }{2}\right )=


v(\xi ,\eta )= \\ =\varphi(\xi )-


\int_0^{\xi } \varphi(s)\frac {\partial }


{\partial s}J_0[\sqrt {\lambda \eta (\xi-s)}]ds+


\psi (\eta )-\int_0^{\eta } \psi (t)\frac {\partial }


{\partial t}J_0[\sqrt {\lambda \xi (\eta -t)}]dt.


\end{multline*}


Отсюда находим


\begin{multline*}


\frac {\partial ^2v}


{\partial \xi \partial \eta}=-\frac {\lambda }


{4}\varphi(\xi )-


\int_0^{\xi } \varphi(s)\frac {\partial }


{\partial s}\frac {\partial ^2}


{\partial \xi \partial \eta }J_0[\sqrt {\lambda \eta (\xi-s)}]ds-


\frac {\lambda }


{4}\psi (\eta )-\int_0^{\eta } \psi (t)\frac {\partial }


{\partial t}\frac {\partial ^2}


{\partial \xi \partial \eta }J_0[\sqrt {\lambda \xi (\eta -t)}]dt= \\


%


=-\frac {\lambda }{4}\varphi(\xi )+


\frac {\lambda }{4}\int_0^{\xi } \varphi(s)\frac {\partial }


{\partial s}J_0[\sqrt {\lambda \eta (\xi-s)}]ds-


\frac {\lambda }{4}


\psi (\eta )+\frac {\lambda }


{4}\int_0^{\eta } \psi (t)\frac {\partial }


{\partial t}J_0[\sqrt {\lambda \xi (\eta -t)}]dt=-\frac {\lambda }{4}v.


\quad\square


\end{multline*}





Под регулярным в $D$ решением уравнения (4) будем понимать


функцию $u(x,y)$, удовлетворяющую условиям $u(x,y)\in C^1(D)\cap


C^2(D_+\cup D_-)$, $Lu(x,y)\equiv 0$ в $D_+\cup D_-$.





\begin{thm}[{\series{m}\shape{n}\selectfont [2]}]


Если функция $u_0(x,y)$ в области $D$


является регулярным решением уравнения


\begin{equation}


u_{0xx}+\sgn y\cdot u_{0yy}=0,


\end{equation}


то


\begin{equation}


u(x,y)=u_0(x,y)-\int_0^1 u_0(xt,yt)\frac {\partial }


{\partial t}J_0[\sqrt {\lambda (x^2+\sgn y\cdot y^2)(1-t)}]dt


\end{equation}


является в области $D$ регулярным решением уравнения $(4)$.


\end{thm}





{\bf Доказательство} этой теоремы проводится аналогично [2] c использованием


леммы 1 и результатов работы [1].


Отметим, что данная теорема доказана в [2] при


условии $u(0,0)=0$.





Заметим, что уравнение (8) однозначно обратимо относительно функции


$u_0(x,y)$ в классе функций $C(\overline D)$. Действительно,


равенство (8) перепишем в виде


\begin{equation}


\wt u(r)=\wt u_0(r)-\int_0^r\wt u_0(s) \frac {\partial }{\partial s}


J_0[\sqrt {\lambda r(r-s)}]ds, \tag{8$_1$}


\end{equation}


где


$$


r^2=x^2+y^2\sgn y,\quad u(x,y)=u\bigg(\frac {x}{r}r,\frac {y}{r}r\bigg)=


\wt u(r),\quad u_0\bigg(\frac {x}{r}s,\frac {y}{r}s\bigg)=\wt u_0(s).


$$


Тогда в силу результатов [1], [3] решением уравнения $(8_1)$


является функция вида


\begin{equation}


\wt u_0(r)=\wt u(r)+\int_0^r \wt u(s)\frac {r}{s}


\frac {\partial }{\partial r}


I_0[\sqrt {\lambda s(r-s)}]ds, \tag{8$_2$}


\end{equation}


где $I_0(\cdot )$ -- модифицированная функция Бесселя.





Если функции $\wt u_0(r)$ и $\wt u(r)$


непрерывны в $\overline D$, то равенства $(8_1)$ и $(8_2)$


являются формулами взаимного обращения [3].





Таким образом, справедлива 





\begin{thm}


Если функции $u_0(x,y)$ и $u(x,y)$


непрерывны в $\overline D$ и являются соответственно регулярными


в $D$ решениями уравнений $(7)$ и $(4)$, то между ними


существует взаимнооднозначное соответствие, которое


устанавливается по формулам $(8_1)$ и $(8_2)$.


\end{thm}





Pассмотрим в области $D$ задачу Трикоми для уравнения (4).





{\it Задача T.}


Найти в области $D$ функцию $u(x,y)$,


удовлетворяющую условиям


\begin{gather}


u(x,y) \in C(\overline D)\cap C^1(D) \cap C^2(D_+\cup D_-);\\


%


Lu(x,y)\equiv 0 \quad \text{при}\quad (x,y)\in D_+\cup D_-;\\


%


u(x,y)\big|_\Gamma=u(x(s),y(s))=\varphi (s),


\quad 0\le s\le l,


\end{gather}


$x=x(s)$, $y=y(s)$ --- параметрические уравнения 


кривой $\Gamma$, $s$ --- длина дуги, отсчитываемая


от точки $B$, $l$ --- длина кривой $\Gamma;$


\begin{equation}


u(x,y)\big|_{AC}=u(x,-x)=\psi (x),\quad 0\le x \le 1/2,


\end{equation}


где $\varphi (l)=\psi (0)$, $\varphi (s)$ и $\psi (x)$ ---


заданные достаточно гладкие функции.


\medskip





Интегральное представление (8) позволяет свести решение


задачи Трикоми для уравнения (4) в области $D$ с краевыми данными


$u=\varphi $ на $\Gamma$ и $u=\psi$ на $AC$ к решению задачи


Трикоми для уравнения (7) в области $D$ с краевыми условиями


$u_0=\varphi_0$ на $\Gamma$ и $u_0=\psi_0$ на $AC$.





Действительно,


прежде всего заметим, что $u(x,y)=u_0(x,y)$ на $AC(x+y=0)$, поэтому


$\psi_0(x)=\psi(x)$.


Для нахождения функции $\varphi_0(s)$ в области $D_+$ функцию $u_0(x,y)$


определим


как решение задачи Хольмгрена для уравнения Лапласа с граничными условиями


\begin{gather*}


u_0(x,y)\big|_\Gamma=u_0(x(s),y(s))=


\varphi_0(s),\quad 0\le s\le l, \\


%


\frac {\partial u_0(x,y)}{\partial y}\bigg|_{y=0}


=\nu_0(x),\quad 0<x<1.


\end{gather*}


Методом Грина решение этой задачи


выписывается в явном виде ([4], гл.\,4; [5])


\begin{equation}


u_0(x,y)=\int_0^1\nu_0(\xi )G(\xi,0;x,y)d\xi +


\int_0^l\varphi_0(s)


\frac {\partial G(\xi (s),\eta (s);x,y)}{\partial N}ds,


\end{equation}


где


$G(\xi ,\eta ;x,y)$


--- функция Грина задачи Хольмгрена, а функция $\nu_0(x)$ определяется


соотношением


\begin{equation}


\nu_0(x)=F(x)+\int_0^1F(t)R(t,x)dt,


\end{equation}


где $R(t,x)$ --- резольвента известного ядра интегрального уравнения


типа Фредгольма,


\begin{equation}


F(x)=\frac {1}{2\pi }\int_0^l \varphi_0(s)\frac {\partial }


{\partial x}\frac {\partial }


{\partial N}G[\xi (s),\eta (s);x,0]ds-


\frac {1}{2}\psi '_0\bigg(\frac {x}{2}\bigg).


\end{equation}


Как видим, функция $\nu_0(x)$ определяется через функции


$\varphi_0(s)$ и $\psi_0(x)$. Но функция $\psi_0(x)$


известна и равна $\psi (x)$.





Подставляя (15) в (14), а последнее --- в (13), получим


\begin{gather}


u_0(x,y)=\int_\Gamma \varphi_0(s)K(s;x,y)ds+f(x,y),\\


%


K(s;x,y)=\frac {1}{2\pi }\int_0^1G(\theta,0;x,y) \frac {\partial }


{\partial \theta}\frac {\partial }


{\partial N}G[\xi (s),\eta (s);\theta ,0]d\theta + \notag \\


%


+\int_0^1 G(\theta ,0;x,y)\int_0^1R(t,\theta)\frac {\partial }


{\partial t}\frac {\partial }


{\partial N}G[\xi (s),\eta (s);t,0]dt\,d\theta, \notag \\


%


f(x,y)=-\frac {1}{2}\int_0^1\psi '_0\bigg(\frac{\theta }{2}\bigg)


G(\theta ,0;x,y)d\theta -\frac {1}{4\pi }


\int_0^1 G(\theta ,0;x,y)\int_0^1\psi '_0


\bigg(\frac{t}{2}\bigg)R\,dt\,d\theta.\notag


\end{gather}


Функцию $u_0(x,y)$, заданную формулой (16), подставим в интегральное


слагаемое формулы (8). Тогда имеем


\begin{multline}


u(x,y)=u_0(x,y)-\int_0^l \varphi_0(s)\int_0^1 K(s;xt,yt)\frac {\partial }


{\partial t}J_0[\sqrt {\lambda (x^2+y^2)(1-t)}]dt\,ds-\\


%


-\int_0^1 f(xt,yt)\frac {\partial }


{\partial t}J_0[\sqrt {\lambda (x^2+y^2)(1-t)}]dt.


\end{multline}


Переходя в (17) к пределу при $(x,y)\to (x(s),y(s))\in \Gamma$,


получим интегральное уравнение относительно функции


\begin{equation}


\varphi_0(s)-\int_0^l\varphi_0(\tau )H(\tau ,s)d\tau= g(s),


\end{equation}


где


\begin{align*}


H(\tau ,s)&=\int_0^1 K(\tau ;x(s)t,y(s)t)\frac {\partial }


{\partial t}J_0[\sqrt {\lambda [x^2(s)+y^2(s)](1-t)}]dt,\\


%


g(s)&=\varphi (s)+\int_0^1 f[x(s)t,y(s)t]\frac {\partial }


{\partial t}J_0[\sqrt {\lambda [x^2(s)+y^2(s)](1-t)}]dt.


\end{align*}





Если функция $\varphi (s)$ непрерывна на $[0,l]$ и


в достаточно малой окрестности


точек $s=0$ и $s=l$ удовлетворяет условию Г\"ельдера с показателем


$\alpha \in [\frac {1}{2},1]$, $\psi (x)\in C^1[0,\frac {1}{2}]


\cap C^2(0,\frac {1}{2})$, $\psi ''(x)\in L_2[0,1]$, $\varphi (0)=


\varphi (l)=\psi (0)=0$,


то (18) является интегральным уравнением Фредгольма второго рода.


Разрешимость его вытекает из теоремы 2 и следующих утверждений 


единственности


решения задачи Трикоми для уравнения (4).





\begin{thm}


Если существует решение задачи


$(9)$--$(12)$ для уравнения $(4)$,


то оно единственно при всех комплексных


$\lambda$, удовлетворяющих условиям


\begin{alignat}{2}


|\im \lambda |&< \sqrt {2 }\left (\frac {\pi }{t_2-t_1}\right )^2 +


\sqrt {2}\re \lambda


&\quad \text{при} \quad \re \lambda &<0;\\


|\im \lambda|&<\sqrt {2}\,\frac {9}{4}\left (\frac {\pi}{t_2-t_1}\right )^2


&\quad \text{при} \quad \re\lambda &\ge 0,


\end{alignat}


где


$t_2=\max\limits_{\overline D} 2y$, $t_1=\min\limits_{\overline D} 2y$.


\end{thm}





\begin{pf}


Пусть $\re u(x,y)=u_1$, $\im u(x,y)=u_2$,


$\lambda =\re\lambda +i \im \lambda$. Тогда задача (9)--(12)


равносильна задаче Трикоми для системы


$$


U_{xx}+\sgn y\cdot U_{yy}+CU=0,


$$


где


$$


C=\begin{pmatrix}


-\re \lambda & \im \lambda\\


- \im \lambda & - \re \lambda


\end{pmatrix},\quad U=(u_1,u_2).


$$





Теперь, применяя принцип максимума модуля решений систем


уравнений смешанного типа [6], получим условия (19) и (20)


относительно параметра $\lambda$. При выполнении условий


(19), (20) и $u=0$ на $AC$ следует,


что $\max\limits_{\overline D}|U|$ достигается


на кривой $\Gamma$. Отсюда уже следует единственность решения задачи (9)--(12)


для уравнения (4). Более полное доказательство этой


теоремы приведено в ([7], гл.\,3, \S\,5).


\end{pf}





\begin{thm}


Если существует решение задачи $(9)$--$(12)$ для


уравнения $(4)$, то оно единственно при всех комплексных $\lambda$,


удовлетворяющих условию 


\begin{equation}


|\lambda |<2\alpha_0-\re \lambda ,\quad \alpha_0=\frac {4}{9\mes D_+} .


\end{equation}


\end{thm}





{\bf Доказательство} этой теоремы опирается на следующую лемму.





\begin{lem}


Если $u(x,y)=0$ на $AC$, $a=|\mu_2|$,


то для любого регулярного в области $D$ решения уравнения


$(4)$ имеет место равенство


\begin{equation}


\re\int_0^1\overline u 


\frac{\partial u}{\partial y}\bigg|_{y=0}


e^{-2ax}dx\ge 0,


\end{equation}


где $\mu ^2=\lambda$, $\mu =\mu_1+i\mu_2$.


\end{lem}





\begin{pf}


Предварительно для уравнения (4) в области $D_-$


построим решение задачи Дарбу с краевыми условиями


$u(x,0)=\tau (x)$, $0\le x\le 1$, $u(x,-x)=0$, $0\le x\le \frac {1}{2}$,


$\tau (0)=0$.





Решение этой задачи построено [8] в явном виде


\begin{equation}


u(x,y)=\tau (x+y)+y\sqrt {\lambda }\int_0^{x+y}


\frac {J_1\sqrt {\lambda (x+y-t)(x-y-t)}\tau (t)}{\sqrt{(x+y-t)(x-y-t)}}dt,


\end{equation}


где $J_1(\cdot )$ --- функция Бесселя первого рода, $\sqrt {\lambda }>0$


при $\lambda >0$. Используя (23), вычислим


\begin{equation}


u_y(x,0-0)=\nu (x)=\tau '(x)+\mu \int_0^x\tau (t)


\frac {J_1[\mu (x-t)]}{(x-t)}dt.


\end{equation}


Продолжив функцию $\tau (x)=0$ при $x\ge 1$,


равенству (24) придадим несколько иной вид. Bоспользовавшись


формулой преобразования Лапласа свертки ([9], с.\,504) для


$$


\int_0^x \mu \tau (x)\frac {J_1[\mu (x-t)]}{x-t}dt,


$$


получим


\begin{equation}


\nu (x)=\tau '(x)+\frac {1}{2\pi i}\int_{a-i\infty }^{a


+i\infty }(\sqrt {p^2+\mu ^2}-p)\overline {\tau }(p) \exp(px)dp,


\end{equation}


где $\re p\ge |\im \mu |$, $a\ge |\im \mu |$,


$\tau (x)=\tau_1(x)+i\tau_2(x)$,


$\overline {\tau }(x)=\tau_1(x)-i\tau_2(x)$.


Тогда аналогично ([10], c.\,15) на основании равенства (25) рассмотрим


интеграл


\begin{multline}


I=\int_0^1\overline {\tau }(x)\nu (x)e^{-2ax}dx=


\int_0^1\overline {\tau }(x)\tau '(x)e^{-2ax}dx+ \\


%


+\frac {1}{2\pi i}\int_0^1\overline {\tau }(x)


\int_{a-i\infty }^{a+i\infty }(\sqrt {p^2+\mu ^2}-p)


\overline {\tau }(p) e^{px}dp\,e^{-2ax}dx=I_1+I_2.


\end{multline}


B силу того, что $\tau (0)=\tau (1)=0$,


\begin{equation}


I_1=-\int_0^1\overline \tau{}'(x)\tau (x) e^{-2ax}dx 


+2a\int_0^1|\tau (x)|^2e^{-2ax}dx.


\end{equation}


Из этого равенства следует


$$


\re I_1=a\int_0^1|\tau (x)|^2e^{-2ax}dx.


$$


Используя равенство Парсеваля для преобразования Лапласа


и тот факт, что $\tau (x)=0$ при $x\ge 1$,


последний интеграл перепишем в виде


\begin{equation}


\re I_1=a\int_0^\infty |\tau (x)|^2e^{-2ax}dx=\frac {a}{2\pi i}


\int_{a-i\infty }^{a+i\infty } |\overline {\tau }(p)|^2dp.


\end{equation}


Далее, снова на основании равенства Парсеваля для преобразования


Лапласа интеграл $I_2$ преобразуем к виду


\begin{equation}


\negthickspace I_2=\int_0^{\infty }\overline {\tau }(x)


\frac {1}{2\pi i}\int_{a-i\infty }^{a+i\infty}


(\sqrt {p^2{+} \mu ^2}{-}p)\wt {\tau }(p)e^{px}dp\,e^{-2ax}dx= 


\frac {1}{2\pi i}\int_{a-i\infty }^{a+i\infty }


|\overline {\tau }|^2


(\sqrt {p^2{+}\mu ^2}{-}p)dp.


\end{equation}


Подставляя (27) и (29) в равенство (26), с учетом (28) получим


\begin{multline*}


\re I=\re I_1+\re I_2=\re\frac {1}{2\pi i}\int_{a-i\infty }^{a+i\infty }


|\overline {\tau }(p)|^2(\sqrt {p^2+\mu ^2}-p+a)dp=\\


%


=\frac {1}{2\pi i}\int_{a-i\infty }^{a+i\infty }


|\overline {\tau }(p)|^2\re(\sqrt {p^2+\mu ^2}-p+a)dp.


\end{multline*}


Отсюда $\re I\ge 0$, если $\re(\sqrt {p^2+\mu ^2}-p+a)\ge 0$.


Последнее выражение имеет место тогда, когда $\re\sqrt {p^2+\mu ^2}\ge 0$,


а это возможно при $a=|\mu_2|$.


\end{pf}





Теперь перейдем непосредственно к доказательству теоремы 4.


B области $D_+$ рассмотрим уравнение (4) и


введем вспомогательную функцию $u=v\exp ax$.


Тогда уравнение (4) при $y>0$ преобразуется к виду


$$


Mv=v_{xx}+v_{yy}+2av_x+(a^2+\lambda )v=0.


$$


Рассмотрим в области $D_+$ тождество $\overline v Mv=0$. Интегрируя 


его и применяя формулу Грина, получим 


\begin{multline*}


\int_{D_+}\overline {v}Mv\,dx\,dy=\int_{\partial D_+}


(-\overline {v}v_ydx+\overline {v}v_xdy)-\int_{D_+}


(\overline {v}_xv_x+\overline {v}_yv_y)dx\,dy+\\


%


+2a\int_{D_+}


\overline {v}v_x dx\,dy+(a^2+\lambda )\int_{D_+}|v|^2dx\,dy=0.


\end{multline*}


Отсюда с учетом граничного условия $v\big|_\Gamma=0$ имеем


$$


-\int_{AB}\overline {v}v_{y|y=0}dx-


\int_{D_+}|\nabla v|^2dx\,dy+


2a\int_{D_+}


\overline {v}v_x dx\,dy+(a^2+\lambda )\int_{D_+}|v|^2dx\,dy=0.


$$


У полученного равенства выделим реальную часть


\begin{equation}


\re\int_{AB}\overline {v}v_{y|y=0}dx+


\int_{D_+}|\nabla v|^2dxdy-2a\re\int_{D_+}


\overline {v}v_x dx\,dy-\re(a^2+\lambda )\int_{D_+}|v|^2dx\,dy=0.


\end{equation}


Поскольку 


$$


2\re\int_{D_+}


\overline {v}v_x dx\,dy=\int_{D_+}


\frac {\partial }{\partial x}|v|^2dx\,dy=


\int_{\partial D_+}|v|^2dy=\int_{AB}|v|^2dy=0,


$$


то равенство (30) примет вид 


\begin{equation}


\int_{D_+}|\nabla v|^2dx\,dy-


\re(a^2+\lambda )\int_{D_+}|v|^2dx\,dy+


\re\int_{AB}\overline vv_{y|y=0}dx=0.


\end{equation}


Если $\re(a^2+\lambda )\le 0,$ то из равенства (31) на основании леммы 2


следует, что $|\nabla v|\equiv 0$ в $D_+$, поэтому $u(x,y)\equiv 0$ в 


$\overline D_+$. Тогда 


в силу единственности решения задачи Коши или Дарбу для уравнения (4)


при $y<0$ вытекает $u(x,y)\equiv 0$ в $\overline D_-$.





Если $\re(a^2+\lambda )\ge 0$, то из равенства (31) получим


\begin{equation}


\int_{D_+}|\nabla v|^2dx\,dy\le


(a^2+\re \lambda )\int_{D_+}|v|^2dx\,dy.


\end{equation}


C другой стороны, в силу аналога неравенства Пуанкаре--Фридрихса


([11], c.\,307; [12], c.\,71)


\begin{equation}


\int_{D_+} |v|^2dx\,dy\le 


\frac {9}{4}\mes D_+\int_{D_+} |\nabla v|^2dx\,dy.


\end{equation}


Тогда из (32) и (33) будем иметь


$$


\int_{D_+}|\nabla v|^2dx\,dy\le


(a^2+\re \lambda )\frac {9}{4}\mes D_+\int_{D_+} |\nabla v|^2dx\,dy.


$$


Отсюда $v(x,y)\equiv 0$ в $\overline D_+$ 


при выполнении условия (21).


Cледовательно, $u(x,y)\equiv 0$ в $\overline D$. 





Заметим, что теоремы 3 и 4 верны без каких-либо ограничений 


геометрического характера на кривую $\Gamma$. 


Из доказанных теорем 1--4 следует справедливость следующего утверждения.





\begin{thm}


Если функция $\varphi (s)$ непрерывна на $[0,l]$ и


в достаточно малой окрестности


точек $s=0$ и $s=l$ удовлетворяет условию Г\"ельдера с показателем


$\alpha \in [1/2,1]$, $\psi (x)\in C^1[0,1/2]


\cap C^2(0,1/2)$, $\psi ''(x)\in L_2[0,1/2]$, $\varphi (0)=


\varphi (l)=\psi (0)=0$ и коэффициент $\lambda $ удовлетворяет 


условиям теоремы $3$ или $4$,


то существует единственное решение задачи $(9)$--$(12)$.


\end{thm}





Формулу (8) можно использовать и при решении


задачи Геллерстедта для уравнения (4) с граничными данными при $y<0$ 


на кусках характеристик $x\pm y=0$ для


уравнения (4) и задачи Трикоми с двумя линиями изменения типа для уравнения


$$


\sgn x\cdot u_{xx}+\sgn y\cdot u_{yy}+\lambda u=0.


$$ 





{\bf 2.} Рассмотрим уравнение (4) в области $Н$, ограниченной отрезком


$AA_1$ оси $x=0$, $|y|<1$, характеристикой $A'B$ уравнения (4),


$A'=(0,-1)$, $B=(1,0)$ и кривой $\Gamma$ из класса Ляпунова с концами


в точках $B$ и $A=(0,1)$, лежащей в первой четверти.


Пусть $H_0=H\cap \{y>0\}$, $H_1=H\cap \{-x<y<0,x>0\}$,


$H_2=H\cap\{y<-x\}$, $O=(0,0)$, $P=(^1/_2,-^1/_2)$,


$l$ --- длина кривой~$\Gamma$.


\medskip





{\it Задача Франкля.}


Найти функцию $u(x,y)$, удовлетворяющую условиям


\begin{gather}


u(x,y) \in C(\overline H)\cap C^1(H\cup OA\cup OA'\setminus OP)


\cap C^2(H_0\cup H_1\cup H_2);\\


%


Lu(x,y)\equiv 0, \quad (x,y)\in H_0\cup H_1\cup H_2;\\


%


u(x,y)\big|_\Gamma=u(x(s),y(s))=\varphi (s),


\ \ 0\le s\le l;\\


%


u(0,y)-u(0,-y)=f(y),\quad 0\le y\le 1;\\


%


u_x(0,y)=0,\quad 0<|y|<1,


\end{gather}


где $\varphi (s)$ и $f(y)$ ---


заданные достаточно гладкие функции.





\begin{lem}


Если функция $u_0(x,y)$


является регулярным в области $H$ решением уравнения $(7)$, то функция


\begin{equation}


u(x,y)=\begin{cases}


u_0(x,y)-\int\limits_0^1 u_0(xt,yt)\frac {\partial }


{\partial t}J_0[\sqrt {\lambda (x^2+y^2)(1-t)}]dt, & (x,y)\in H_0;\\[2mm]


%


u_0(x,y)-\int \limits_0^1 u_0(xt,yt)\frac {\partial }


{\partial t}J_0[\sqrt {\lambda |x^2-y^2|(1-t)}]dt, & (x,y)\in H_1\cup H_2,


\end{cases}


\end{equation}


является регулярным в $H$ решением уравнения $(4)$.


\end{lem}





Cправедливость этой леммы следует из леммы 1 и теоремы 1.





Формула (39) позволяет свести решение задачи (34)--(38) для уравнения (4)


в области $Н$ c краевыми условиями


$u=\varphi$ на $\Gamma$ и $u(0,y)-u(0,-y)=f(y)$, $0\le y\le 1$,


к решению задачи Франкля для уравнения (7) в области $Н$


с краевыми условиями


$u_0=\varphi_0$ на $\Gamma$ и $u_0(0,y)-u_0(0,-y)=f_0(y)$, $0\le y\le 1$.





Действительно, пусть $u_0(x,y)$ --- решение задачи Франкля для


уравнения (7) c отмеченными выше граничными условиями. Тогда из


формулы (39) имеем


\begin{alignat*}{2}


u(0,y)&=u_0(0,y)-\int_0^1 u_0(0,yt)\frac {\partial }


{\partial t}J_0[y\sqrt {\lambda (1-t)}]dt, &\quad y&\ge 0;\\


%


u(0,y)&=u_0(0,y)-\int_0^1 u_0(0,yt)\frac {\partial }


{\partial t}J_0[-y\sqrt {\lambda (1-t)}]dt, &\quad y&\le 0,


\end{alignat*}


или


$$


u(0,-y)=u_0(0,-y)-\int_0^1 u_0(0,-yt)\frac {\partial }


{\partial t}J_0[y\sqrt {\lambda (1-t)}]dt,\quad y\ge 0.


$$


Отсюда с учетом условия (37) получим


\begin{multline*}


f(y)=u(0,y)-u(0,-y)=u_0(0,y)-u_0(0,-y)- \\


%


-\int_0^1 [u_0(0,yt)-u_0(0,-yt)]\frac {\partial }


{\partial t}J_0[y\sqrt {\lambda (1-t)}]dt= \\


%


=f_0(y)-\int_0^1 f_0(yt)\frac {\partial }


{\partial t}J_0[y\sqrt {\lambda (1-t)}]dt=


f_0(y)-\int _0^y f_0(s)\frac {\partial }


{\partial s}J_0[\sqrt {\lambda y(y-s)}]ds.


\end{multline*}


Обращая последнее интегральное уравнение по теореме 1 [3],


найдем функцию


$$


f_0(y)=f(y)+\int_0^y f(s)\frac {y}{s}\frac {\partial }


{\partial y}I_0[\sqrt {\lambda s(y-s)}]ds,


$$


которая имеет ту же гладкость, что и функция $f(y)$.





Для нахождения функции $\varphi_0(s)$ поступим так. B


области $H_0$ выпишем методом Грина решение


$u_0(x,y)$ задачи Хольмгрена для уравнения Лапласа аналогично


тому, как было показано в п.\,1, и подставим его в интегральный


член формулы (39) при $y>0$. Затем, переходя к пределу при


$(x,y)\to (x(s),y(s))\in \Gamma$, получим интегральное уравнение


относительно $\varphi_0(s)$, к которому при известных условиях


на функции $\varphi (s)$ и $f(y)$ ([4], с.\,340; [13], гл.\,8) применима


теория Фредгольма.


\medskip





{\bf 3.} Рассмотрим телеграфное уравнение (3) в области $D_-^k$,


ограниченной отрезками $AB$ $(y=0)$, $AC$ $(y=-kx$, $0<k\le 1)$ и


$CB$ $(x-y=1)$, и поставим задачи Дарбу.


\medskip





{\it Задача $D_1$.} 


Найти функцию $u(x,y)$, удовлетворяющую условиям


\begin{gather}


u(x,y) \in C(\overline D{}_-^k)\cap C^2(D_-^k);\\


%


u_{xx}-u_{yy}+\lambda u\equiv 0, \quad (x,y)\in D_-^k;\\


%


u(x,y)\big|_{y=-kx}=\psi (x),


\quad 0\le x\le \frac{1}{1+k};\\


%


u(x,y)\big|_{y=0}=\tau (x),\quad 0\le x\le 1.


\end{gather}





{\it Задача $D_2$.} 


Найти функцию $u(x,y)$,


удовлетворяющую условиям $(40)$--$(42)$ и, кроме того,


$$


u_y(x,0)=\nu (x),\quad 0<x<1,


$$


где $\psi (x)$, $\tau (x)$, $\nu (x)$ --- заданные


достаточно гладкие функции.


\medskip





Обозначим через $u_0(x,y)$ решения задач $D_1$ и $D_2$ для


уравнения (41) при $\lambda =0$ и через


$\psi_0 (x)$, $\tau_0 (x)$, $\nu_0 (x)$ --- cоответствующие граничные 


функции. Тогда справедливa





\begin{lem}


Если $\tau_0(x)\in C[0,1]\cap C^1[0,1)\cap C^2(0,1)$,


$\psi_0(x)\in C[0,\frac {1}{1+k}]\cap C^1[0,\frac {1}{1+k})\cap 


C^2(0,\frac {1}{1+k})$, $\tau_0(0)=\psi_0(0)=0$, функции


$\tau''_0(x\alpha ^n)\alpha ^n$ и $\psi''_0(x\alpha ^n)\alpha ^n$


ограничены по $n$ при любом фиксированном $x$, то функция


\begin{multline}


u_0(x,y)=\tau_0(x+y)-\sum_{n=0}^{\infty}


\bigg\{\psi_0\bigg[(1+\alpha)\alpha ^n\frac {x+y}{2}\bigg]- \\


%


-\psi_0\bigg[(1+\alpha)\alpha ^n\frac {x-y}{2}\bigg]-


\tau_0 [\alpha ^{n+1}(x+y)]+


\tau_0 [\alpha ^{n+1}(x-y)]\bigg\},


\end{multline}


где $\alpha =(1-k)/(1+k)$, является решением задачи $(40)$--$(43)$


при $\lambda =0$. 





Если $\nu_0(x)\in C[0,1)\cap C^1(0,1)\cap L_1[0,1]$,


функция $\nu'_0(x\alpha ^n)\alpha ^n$ ограничена по $n$ при 


любом фиксированном $x$,


а функция $\psi_0(x)$ удовлетворяет отмеченным выше условиям,


то функция


\begin{multline}


u_0(x,y)=\int_0^{x+y}\nu_0(t)dt+\sum_{n=0}^{+\infty}(-1)^n


\bigg\{\psi_0\bigg[(1+\alpha)\alpha ^n\frac {x+y}{2}\bigg]+ \\


%


+\psi_0\bigg[(1+\alpha)\alpha ^n\frac {x-y}{2}\bigg]-


\int_0^{(x+y)\alpha ^{n+1}}\nu_0(t)dt-


\int_0^{(x-y)\alpha ^{n+1}}\nu_0(t)dt\bigg\}


\end{multline}


является решением задачи $D_2$ для уравнения струны. 


\end{lem}





Отметим, что формула (44) при $\psi_0(x)\equiv 0$ приведена


в ([10], c.\,53), и доказательство леммы проводится аналогично 


этой работе.





B силу формул (44) и (45) лемма 1 позволяет решение задач $D_1$ и $D_2$


для уравнения (3) свести к решению задач $D_1$ и $D_2$


для уравнения струны.





Действительно, пусть $u_0(x,y)$ --- решение задачи $D_1$ в области


$D_-^k$ для уравнения струны (5). Тогда функция $u_0(x,y)$ выражается формулой


(44). Подставляя $u_0(x,y)$ в интегральный член формулы (6), получим


\begin{multline}


u(x,y)=u_0(x,y)-\int_0^1 \bigg\{\tau_0(xt+yt)-\sum_{n=0}^{\infty}


\bigg\{\psi_0\bigg[(1+\alpha)\alpha ^nt\frac {x+y}{2}\bigg]-\\


%


-\psi_0\bigg[(1+\alpha)\alpha ^nt\frac {x-y}{2}\bigg]-


\tau_0 [\alpha ^{n+1}t(x+y)]+ 


\tau_0 [\alpha ^{n+1}t(x-y)]\bigg\}\bigg\}


\frac {\partial }


{\partial t}J_0[\sqrt {\lambda (x^2-y^2)(1-t)}]dt.


\end{multline}


Полагая в (46) $y=-kx$, будем иметь


$$


\psi (x)=\psi_0(x)-\int_0^1 \psi_0(xt)\frac {\partial }


{\partial t}J_0[\sqrt {\lambda (1-k^2)x^2(1-t)}]dt


$$


или


\begin{equation}


\psi (x)=\psi_0(x)-\int_0^x \psi_0(s)\frac {\partial }


{\partial s}J_0[\sqrt {\lambda (1-k^2)x(x-s)}]ds.


\end{equation}


Обращая интегральное уравнение (47) по теореме 1 [3],


находим неизвестную функцию 


\begin{equation}


\psi_0(x)=\psi(x)+\int_0^x \psi(t)\frac {x}{t}\frac {\partial }


{\partial t}I_0[\sqrt {\lambda (1-k^2)t(x-t)}]dt,


\end{equation}


которая обладает той же гладкостью, что и функция $\psi (x)$.





Из формулы (46) при $y=0$ получим интегральное уравнение относительно


функции $\tau_0(x)$


$$


\tau (x)=\tau_0(x)-\int_0^1 \tau_0(xt)\frac {\partial }


{\partial t}J_0[\sqrt {\lambda x^2(1-t)}]dt,


$$


решением которого является функция


\begin{equation}


\tau_0 (x)=\tau (x)+\int_0^x \tau(t)\frac {x}


{t}\frac {\partial }{\partial x}I_0[\sqrt {\lambda t(x-t)}]dt.


\end{equation}





Итак, если функции $\psi (x)$ и $\tau (x)$ обладают отмеченными


в лемме 4 свойствами функций $\psi_0(x)$ и $\tau_0(x)$,


которые определены формулами (48) и (49),


то $\psi_0(x)$ и $\tau_0(x)$


удовлетворяют условиям леммы 4. Тогда формула (6), где $u_0(x,y)$ задана


формулой (44), а функции $\psi_0(x)$ и $\tau_0(x)$ определены соответственно


формулами (48) и (49), определяет решение задачи (40)--(43).





B случае задачи $D_2$ аналогично находится функция $\psi_0(x)$ по формуле


(48), а $\nu_0(x)$ определяется как решение интегрального уравнения


$$


\nu (x)=\nu_0(x)+\int_0^1 \nu_0(xt)t\frac {\partial }


{\partial t}J_0[\sqrt {\lambda x^2(1-t)}]dt,


$$


к которому также применима теорема 1 [3].





Поэтому формула (6), в которой $u_0(x,y)$ определена формулой (45),


$\psi_0(x)$ задана равенством (48), функция $\nu_0(x)$ находится


по формуле


$$


\nu_0(x)=\nu (x)+\int_0^x \nu(s)\frac {\partial }


{\partial x}I_0[\sqrt {\lambda s(x-s)}]ds,


$$


определяет решение задачи $D_2$ для телеграфного уравнения (3)


в области $D_-^k$.





Отметим, что на основании этих результатов аналогично п.\,1 можно свести


обобщенную задачу Трикоми для уравнения (4) 


в области $D$, где область гиперболичности совпадает с $D_-^k$,


с граничными условиями $u=\varphi $ на $\Gamma$ и $u=\psi $ на 


$AC(y=-kx,0<k\le 1)$ к обобщенной задаче Трикоми для


уравнения (7) в той же области с краевыми условиями 


$u_0=\varphi_0$ на $\Gamma$ и $u_0=\psi_0$ на $AC$.





\begin{thebibliography}{99}





\bibitem{1}


Bекуа И.Н. {\em Новые методы решения эллиптических уравнений}. -- М.--Л:


Гостехиздат, 1948. -- 296 с.


\bibitem{2}


Жегалов B.И. {\em Об одном случае задачи Трикоми} // Тр. семин.


по краев. задачам. -- Изд-во Казанск. ун-та, 1966. -- Bып.\,3. -- C.\,28--36.


\bibitem{3}


Cабитов К.Б. {\em Обращение некоторых интегральных уравнений Bольтерра}


/ /ДAН CCCР. -- 1990. -- Т.\,314. -- \No\,2. -- C.\,300--303.


\bibitem{4}


Бицадзе A.B. {\em Некоторые классы уравнений в частных производных}. --


М.: Наука, 1981 -- 448~с.


\bibitem{5}


Кальменов Т.Ш. {\em О спектре задачи Трикоми для уравнения 


Лаврентьева--Бицадзе} // Дифференц. уравнения. -- 1977. -- T.\,13. -- 


\No\,8. -- C.\,1718--1725.


\bibitem{6}


Cабитов К.Б. {\em Экстремальные свойства модуля решений одного класса


систем уравнений смешанного типа} // ДAН CCCР. -- 1990. -- T.\,310 -- 


\No\,1. -- C.\,33--36.


\bibitem{7}


Cабитов К.Б. {\em Некоторые вопросы качественной и спектральной теории


уравнений смешанного типа}: Дис. $\dotsc$ докт. физ.-матем. наук.


-- М.: МГУ, 1991. -- 313 с.


\bibitem{8}


Cабитов К.Б. {\em Построение в явном виде решений задач Дарбу для 


телеграфного уравнения и их применение при обращении интегральных уравнений}


// Дифференц. уравнения. -- 1990. -- T.\,26. -- \No\,6. -- C.\,1023--1032.


\bibitem{9}


Лаврентьев М.A., Шабат Б.B. {\em Методы теории функции комплексного 


переменного}. -- М.: Наука, 1973. -- 736 с.


\bibitem{10}


Моисеев Е.И. {\em Уравнения смешанного типа со спектральным параметром}. --


М.: МГУ, 1988. -- 150 с.


\bibitem{11}


Березанский Ю.М. {\em Разложение по собственным функциям самосопряженных 


операторов}. -- Киев.: Наук. думка, 1965. -- 798 c.


\bibitem{12}


Ладыженская О.A., Уральцева Н.Н. {\em Линейные и квазилинейные уравнения 


эллиптического типа}. -- М.: Наука, 1973. -- 576 с.


\bibitem{13}


Cмирнов М.М. {\em Уравнения смешанного типа}.


-- М.: Bысш. школа, 1985. -- 304 с.





\end{thebibliography}





\vskip 2cm





\post{\it Cтерлитамакский государственный}{\it Поступила}


\post{\it педагогический институт}{23.10.2001}


\smallskip


\post{\it Cтерлитамакский филиал}{}


\post{\it Aкадемии наук Республики Башкортостан}{}





\label{end}


\end{document}


